\scanpage{33}
\begin{center}
Cours DEA 1977/78 de A. Grothendieck

\uline{Géométrie arithmétique des polygônes réguliers}
\end{center}

{\reversemarginpar\margin{\handwritten{!}}\normalmarginpar}%
% A short hand-drawn line under the typed words ``en commun''.
Ce ``cours'' se fera dans l'esprit d'une investigation \handwritten{\uline}{en commun} de multiples facettes de la notion de polygône régulier au-dessus d'un corps de base, puis d'un anneau de base quelconque. Ce sera une occasion stimulante pour nous familiariser avec les notions de base de géométrie vectorielle, affine, projective - avec ou sans métrique - en dimensions un, deux et trois, et avec les phénomènes particuliers à ces dimensions, ainsi qu'avec les aspects parfois imprévus que prennent des notions familières (comme celle de réflexion\deleted{s}, ou de centre d'un polygône régulier) sur les corps de caractéristique positive, et notamment sur les corps finis. Le besoin d'une géométrie sur des anneaux de base généraux (tel l'anneau ``sur lequel s'écrit'' le $\mathrm{n}$-polygône régulier ``le plus général'') se fera sentir en cours de route, et nous amènera sans doute à affleurer quelques points de fondements de géométrie algébrique. Si le temps le permet et si d'autres investigations ne nous retiennent , nous tenterons, dans la volée, de construire l'icosaèdre régulier ``général'' sur un anneau de base convenable, fini sur l'anneau \hspace{.5em}${\handwritten{\mathbb}{Z}}$\hspace{.5em} des entiers, ayant des corps résiduels de toutes caractéristiques. (La caractéristique deux \hspace{.5em}donnera sans doute du fil à retordre !).

A partir d'un point de départ intuitif et familier, ce ``cours'' voudrait être une ouverture à un univers fascinant, pleins de mystères qui ne demandent qu'à être élucidés. Il ne peut donner des fruits qu'avec une participation active de l'étudiant et un travail personnel important, notamment pour développer a\deleted{e}vec le détail nécessaire les nombreux points\handwritten{(}de fondement\inserted{\handwritten{s}} notamment) qui n'auront été qu'esquissés dans le ``cours''. Il lui est très fortement conseillé de suivre en même temps le cours \hspace{.5em}C4\hspace{.5em} sur l'icosaèdre, qui se fera dans un esprit analogue. Il y aura sans doute de nombreuses interférences entre ces deux cours de réflexion.
% The underlying e in avec has a visible mid-height horizontal bar.
% The live v is clear; its heavy rounded overstrike leaves the inscription method open.
% The opening parenthesis before de fondements is retained as handwritten at medium confidence;
% reinforcement of an existing type impression remains open.
% Z has two obliques. The added right oblique is assigned to hand at medium confidence;
% an additional type impression remains open as to instrument, not as to the doubled glyph.
% C4: both characters are clear and retained without a glyph doubt mark. The tall narrow C,
% irregular terminal joins, and unfooted descending 4 leave inscription method unresolved.
% Detached geometry, origin and linguistic role unassigned. Source-coordinate positions
% are retained within each TikZ picture; these marks are not added to the verbal text.
% Upper paragraph: pale loop immediately below the a of affine (about x1258,y645),
% and compact pale point below de before réflexion (about x1019,y877).
\figure{\begin{tikzpicture}[x=0.008cm,y=-0.008cm]
\path[use as bounding box] (1005,632) rectangle (1270,888);
\draw[black!35,line width=.55pt] (1254,639) .. controls (1253,643) and (1255,648) .. (1259,648) .. controls (1262,647) and (1262,644) .. (1261,641);
\fill[black!45] (1018.8,877.3) ellipse [x radius=1.2,y radius=1.0];
\end{tikzpicture}}
% Lower paragraph: hook below the end of voudrait, two separated points below active,
% and a connected two-lobed mark at the physical left edge beside le détail.
\figure{\begin{tikzpicture}[x=0.008cm,y=-0.008cm]
\path[use as bounding box] (20,1573) rectangle (1545,1850);
\draw[line width=.75pt,line cap=round] (1500,1579) .. controls (1501,1582) and (1498,1584) .. (1495,1584);
\fill (1534,1723) ellipse [x radius=1.6,y radius=1.2,rotate=-42];
\fill (1528.5,1730) ellipse [x radius=1.3,y radius=1.1];
\fill (27,1841) .. controls (27,1838) and (30,1838) .. (31,1840) .. controls (36,1839) and (38,1841) .. (35,1845) .. controls (32,1847) and (30,1844) .. (30,1843) -- cycle;
\end{tikzpicture}}
% archive sheet number: 32
